\documentclass[../../main.tex]{subfiles}% 注意这里的文件路径不能用 ./main.tex ,否则用latexmk编译子文件会报错
\graphicspath{{\subfix{./image/}}{\subfix{../../image/}}} % 指定图片引用路径,后续可以直接使用图片文件名
% 如果图片存放在上级目录,则使用 ../../image/ 作为备用路径

% 例如:
% \begin{figure}[H]
% \centering
% \includegraphics[scale=0.3]{图.png}
% \caption{}
% \label{figure:图}
% \end{figure}
% 注意:上述\label{}一定要放在\caption{}之后,否则引用图片序号会只会显示??.

\begin{document}

\section{Dupin标形和曲面参数方程在一点的标准展开}

\begin{definition}
设正则曲面$S$在点$p$的主曲率为$\kappa_1,\kappa_2$，对应一对正交主方向单位向量$\boldsymbol{e}_1,\boldsymbol{e}_2$，由Euler公式，曲面沿任意切方向$\boldsymbol{e}=\cos\theta\boldsymbol{e}_1+\sin\theta\boldsymbol{e}_2$的法曲率满足
\begin{align}
\kappa_n=\kappa_1\cos^2\theta+\kappa_2\sin^2\theta.\label{eq:euler_pointtype}
\end{align}
在切平面直角坐标系$\{p;\boldsymbol{e}_1,\boldsymbol{e}_2\}$中，取切方向射线上的点$q$满足$|pq|=\frac{1}{\sqrt{|\kappa_n(\theta)|}}$，其坐标为
\begin{align*}
x=\frac{1}{\sqrt{|\kappa_n(\theta)|}}\cos\theta,\quad y=\frac{1}{\sqrt{|\kappa_n(\theta)|}}\sin\theta,
\end{align*}
则坐标满足方程
\begin{align}
\kappa_1x^2+\kappa_2y^2=\mathrm{sign}(\kappa_n(\theta)).\label{eq:dupin_equation}
\end{align}
称该方程的轨迹为曲面$S$在点$p$的\textbf{Dupin标形}，直观刻画该点处法曲率的变化规律。
\end{definition}

\begin{definition}
设$K=\kappa_1\kappa_2$为曲面在点$p$的Gauss曲率，据此将曲面上的点分为三类：
\begin{enumerate}[(1)]
\item 若$K>0$，称$p$为\textbf{椭圆点}，Dupin标形为椭圆；圆点属于椭圆点；
\item 若$K<0$，称$p$为\textbf{双曲点}，Dupin标形为两对共轭双曲线，$\kappa_n=0$的方向为渐近方向；
\item 若$K=0$，称$p$为\textbf{抛物点}；其中$\kappa_1,\kappa_2$不全为零时Dupin标形为一对平行直线，$\kappa_1=\kappa_2=0$（平点）时无Dupin标形。
\end{enumerate}
平点属于抛物点。
\end{definition}

\begin{figure}[H]
\centering
\includegraphics[scale=0.5]{Dupin标形.png}
\caption{Dupin标形}
\label{figure:dupin_indicatrix}
\end{figure}

\begin{table}[H]
\centering
\caption{点型、Dupin标形和近似曲面}
\label{tab:surf_points}
\begin{tabular}{l c c c c l}
\toprule
\textbf{点型} & \textbf{Gauss曲率} & \textbf{Dupin标形} & \textbf{近似曲面} & \textbf{渐近方向} & \textbf{备注} \\
\midrule
椭圆点      & $K>0$   & 椭圆       & 椭圆抛物面   & 无   & 两对共轭 \\
双曲点      & $K<0$   & 双曲线     & 双曲抛物面   & 两个 & \\
抛物线点 (非平点) & $K=0$ & 一对平行直线 & 抛物柱面     & 一个 & \\
抛物线点 (平点)   & $K=0$ & 任意       & 无           & 无   & 切方向 \\
\bottomrule
\end{tabular}
\end{table}

\begin{definition}
在点$p$附近取适配参数系，满足$E=G=1,\ F=M=0,\ L=\kappa_1,\ N=\kappa_2$，将曲面参数方程在$p$处作Taylor展开，忽略高阶无穷小，得到二次曲面
\begin{align}
z=\frac{1}{2}(\kappa_1x^2+\kappa_2y^2),\label{eq:approximate_surface}
\end{align}
称其为曲面$S$在点$p$的\textbf{近似曲面}$S^*$。
\end{definition}

\begin{proposition}
近似曲面$S^*$与原曲面$S$在点$p$相切，二者具有相同的主方向、主曲率，沿任意切方向的法曲率一致。点型与近似曲面的对应关系为：
\begin{enumerate}[(1)]
\item 椭圆点：$\kappa_1,\kappa_2$同号，近似曲面为椭圆抛物面，曲面在该点附近全部落在切平面的一侧；
\item 双曲点：$\kappa_1,\kappa_2$异号，近似曲面为双曲抛物面，曲面在该点附近跨切平面两侧；
\item 非平点抛物点：一主曲率为零，近似曲面为抛物柱面；
\item 平点：无二次近似曲面。
\end{enumerate}
\end{proposition}

\begin{figure}[H]
\centering
\includegraphics[scale=0.5]{近似曲面.png}
\caption{近似曲面}
\label{figure:approximate_surf}
\end{figure}

\begin{example}
考察环面$S:\boldsymbol{r}=((a+r\cos u)\cos v,(a+r\cos u)\sin v,r\sin u)\ (0<r<a)$上各类点的分布。
\end{example}
\begin{solution}

计算第一、第二基本量：
\begin{align*}
E=r^2,\quad F=0,\quad G=(a+r\cos u)^2,\\
L=r,\quad M=0,\quad N=\cos u(a+r\cos u).
\end{align*}
Gauss曲率为
\begin{align*}
K=\frac{LN-M^2}{EG-F^2}=\frac{\cos u}{r(a+r\cos u)}.
\end{align*}
由$\cos u$的符号判定点型：
\begin{enumerate}[(1)]
\item 当$u=\frac{\pi}{2}$或$\frac{3\pi}{2}$时，$K=0$，为抛物点，分布于环面上下两条平行圆；
\item 当$\frac{\pi}{2}<u<\frac{3\pi}{2}$时，$K<0$，为双曲点，分布于环面内侧；
\item 当$0\leqslant u<\frac{\pi}{2}$或$\frac{3\pi}{2}<u\leqslant2\pi$时，$K>0$，为椭圆点，分布于环面外侧。
\end{enumerate}
椭圆点与双曲点分别构成曲面的开子集，二者的边界由抛物点组成。

\end{solution}

\begin{example}
猴鞍面$z=x^3-3xy^2$，原点为孤立的平点。(见\reffig{figure:平点不构成开集的曲面}(a))
\end{example}

\begin{example}
曲面$z=x^4y^4$，平点分布于$x$轴与$y$轴上。(见\reffig{figure:平点不构成开集的曲面}(b))
\end{example}
\begin{remark}
若曲面上的平点构成开集，则该区域局部为平面；孤立平点或沿曲线分布的平点附近，曲面无统一的二次近似模型，几何形态较为复杂。
\end{remark}

\begin{example}
如果旋转曲面上的经线有切向量与旋转轴垂直,则相应的切点是该曲面的抛物点.
\end{example}

\begin{figure}[H]
\centering
\includegraphics[scale=0.5]{平点不构成开集的曲面.png}
\caption{近似曲面}
\label{figure:平点不构成开集的曲面}
\end{figure}

\begin{example}
求曲面$\bm{r}=(u^3,v^3,u+v)$上的抛物点的轨迹.
\end{example}

\begin{example}
研究习题4.4中第5题的管状曲面上各种类型点的分布情况.
\end{example}

\begin{example}
设$\theta$是曲面$S$在双曲点$p$处的两个渐近方向的夹角,则
\begin{enumerate}[(1)]
\item \begin{align*}
\tan\theta=\frac{\sqrt{-K}}{H};
\end{align*}
\item \begin{align*}
\cos\theta=\pm\frac{EN-2FM+GL}{\sqrt{(EN-GL)^2+4(EM-FL)(GM-FN)}},
\end{align*}
\end{enumerate}
其中$E,F,G,L,M,N$分别是曲面$S$在点$p$处的第一类基本量和第二类基本量.
\end{example}

\begin{example}
如果曲面$S$上的渐近曲线网的夹角是常数,则曲面$S$的Gauss曲率$K$和平均曲率$H$的平方成比例.

\end{example}

\begin{example}
求下列曲面在原点处的近似曲面:
\begin{enumerate}[(1)]
\item $z=\exp(x^2+y^2)-1$;
\item $z=\ln\cos x-\ln\cos y$;
\item $z=(x+3y)^2$.
\end{enumerate}
\end{example}

\begin{example}
求曲面$\ln z=-\frac{x^2+y^2}{2}$的Gauss曲率,画出它的草图,并且指出椭圆点和双曲点分布的区域.
\end{example}

\begin{example}
证明:如果曲面$S$在点$p$有三个两两不共线的渐近方向,则该点必定是曲面$S$的平点.
\end{example}


\end{document}